\section{有界集和囿空间}\label{sec:yosida-01-07}
\noindent\textit{原书位置：中译本第 46--48 页；英文版第 44--46 页。}

\begin{definition}[拓扑有界集与囿空间]
拓扑向量空间中的 $B$ 称为有界集，若每个零邻域 $U$ 都吸收 $B$：存在 $c>0$ 使 $B\subset cU$（可先把 $U$ 换成均衡零邻域）。局部凸空间称为囿空间或 bornological 空间，若每个吸收全部有界集的绝对凸集都是零邻域。
\end{definition}
\begin{proposition}[有界集的半范数判别]\label{prop:y1-bounded}
局部凸空间中，$B$ 有界当且仅当对每个连续半范数 $p$，$\sup_{x\in B}p(x)<\infty$。连续线性映射把有界集映为有界集。紧集以及收敛序列连同其极限均有界。
\end{proposition}
\begin{proof}
有界性应用于 $\{p<1\}$ 得一向；反向对一个含有限个半范数条件的零邻域，选择足够大的共同缩放常数即可。若 $T$ 连续，给定目标零邻域 $V$，取源零邻域 $U$ 使 $T(U)\subset V$，则 $B\subset cU$ 推出 $T(B)\subset cV$。紧集的有界性由零邻域的吸收性和有限子覆盖证明；收敛序列加极限是紧集。这最后一条对一般拓扑向量空间也成立。
\end{proof}

\begin{theorem}[囿空间的等价刻画]\label{thm:y1-bornological}
局部凸空间 $X$ 是囿空间，当且仅当每个在全部有界集上有界的半范数都连续。因而从囿空间到局部凸空间的线性映射连续，当且仅当它把有界集映为有界集。
\end{theorem}
\begin{proof}
设 $X$ 囿，半范数 $p$ 在各有界集上有界。则 $\{p\le1\}$ 吸收所有有界集，故是零邻域；由半范数的缩放得 $p$ 连续。反向设每个这种半范数连续。若绝对凸集 $U$ 吸收全部有界集，它尤其吸收每个单点，因此 $p_U$ 有定义；吸收有界集保证 $p_U$ 在有界集上有界，于是连续，$\{p_U<1\}\subset U$ 证明 $U$ 是邻域。

对线性映射 $T:X\to Y$，若它保有界集，则每个连续目标半范数 $q$ 的复合 $q\circ T$ 在源有界集上有界，故连续；半范数刻画推出 $T$ 连续。反向已由上一命题证明。
\end{proof}
赋范空间都是囿空间：若半范数 $p$ 在范数单位球上有界于 $C$，则 $p(x)\le C\|x\|$，所以连续。不要求赋范空间完备。

\begin{theorem}[拟赋范空间上保有界集的线性映射]
设 $X,Y$ 均为定义 \ref{def:y1-Fnorm} 意义下的拟赋范空间。线性映射 $T:X\to Y$ 连续，当且仅当它保拓扑有界集。
\end{theorem}
\begin{proof}
只证明充分性。若 $x_j\to0$，可选正整数 $n_j\to\infty$，使 $n_jx_j\to0$：对每个整数 $k$，固定数乘的连续性给出 $kx_j\to0$，选择递增的 $N_k$ 使 $j\ge N_k$ 时 $q_X(kx_j)<1/k$，再令 $n_j=k$ 对 $N_k\le j<N_{k+1}$。于是 $\{n_jx_j\}$ 有界，故 $\{n_jTx_j\}$ 有界。拓扑有界性保证有界集乘以趋零标量时一致趋于零，因此 $Tx_j=n_j^{-1}(n_jTx_j)\to0$。度量空间上线性映射的序列连续等价于连续，故结论成立。
\end{proof}
\begin{corollary}[补充：可度量局部凸空间是囿空间]
每个可度量的局部凸空间都是囿空间。
\end{corollary}
\begin{proof}
对在有界集上有界的半范数 $p$，将它看作线性映射 $X\to X/\ker p$ 的目标范数。此映射保有界集，上一定理（目标赋范）保证连续，故 $p$ 连续，再引用定理 \ref{thm:y1-bornological}。
\end{proof}
囿空间关注“有界性足以检测连续性”；它与以后出现的桶空间不同。后者的定义要求所有闭的吸收绝对凸集都是邻域。两者的假设和用途需要分别检查。
